% Step 4 research note: Weyl profile concentration, bubble restriction,
% and the Armstrong--LeBrun tension
% S0910A03 Agent2 (Constructor) — Explore
%
% This note extends the Step 3 near-equality analysis with three new results:
% (A) L² concentration of |W+| at the calibrated value 1/(2√6),
% (B) restriction of noncollapsed smooth blow-up limits to locally hyperkähler,
% (C) identification of a structural contradiction between the profile and
%     the equality-mechanism limit via Armstrong's theorem.

\section{The Weyl profile concentration}

\noindent\textbf{Proposition A} (Weyl profile \(L^2\) concentration).
\emph{Under the hypotheses and notation of the Step~3 near-equality analysis
(critical Ishida--LeBrun family, \(\delta_i\to0\), and the energy profile (3)--(7)),}
\[
  \int_M\Bigl(|W_{+,i}|-\tfrac{1}{2\sqrt6}\Bigr)^2d\mu_i
  =O(A\delta_i^2).                                          \tag{A1}
\]

\begin{proof}
The completion-square identity (Step~3, equation~(4)) gives
\[
  \frac{1}{864\pi^2}\int_M(s_i+8\sqrt6\,|W_{+,i}|)^2\,d\mu_i\to0.
\]
We write \(w_i=|W_{+,i}|\).  Since \(8\sqrt6\cdot\tfrac{1}{2\sqrt6}=4\),
\[
  s_i+8\sqrt6\,w_i
  =(s_i+4)+8\sqrt6\bigl(w_i-\tfrac{1}{2\sqrt6}\bigr).
\]
Squaring, integrating, and using \(|s_i+4|\le4\delta_i\):
\[
  384\int(w_i-\tfrac{1}{2\sqrt6})^2\,d\mu_i
  \le\int(s_i+8\sqrt6\,w_i)^2\,d\mu_i
    +16\delta_i^2\,\operatorname{Vol}_i
    +64\sqrt6\,\delta_i\,\operatorname{Vol}_i^{1/2}
    \Bigl[\int(w_i-\tfrac{1}{2\sqrt6})^2\Bigr]^{1/2}.
\]
Setting \(Y_i=\int(w_i-\tfrac{1}{2\sqrt6})^2\) and using
\(\operatorname{Vol}_i\to32\pi^2A/9=O(A)\), this is a quadratic
\[
  384Y_i-C_1A^{1/2}\delta_i Y_i^{1/2}-C_2A\delta_i^2\le0,
\]
which forces \(Y_i^{1/2}=O(A^{1/2}\delta_i)\), hence \(Y_i=O(A\delta_i^2)\).
\end{proof}

\noindent\textbf{Interpretation.}
The ``calibrated value'' \(1/(2\sqrt6)\approx0.204\) is the unique pointwise
Weyl magnitude compatible with \(s\to-4\) and the completion-square residual
vanishing.  It is exactly \(\tfrac14\) of the K\"ahler--Einstein value
\(|W_+|_{\rm KE}=2/\sqrt6\approx0.816\) for metrics with \(\operatorname{Ric}=-g\).

\section{Exclusion of K\"ahler--Einstein bulk and other known geometries}

\noindent\textbf{Corollary B1} (K\"ahler--Einstein exclusion).
\emph{Any subset \(U_i\subset M_i\) on which \(|W_{+,i}|\to2/\sqrt6\) uniformly
satisfies}
\[
  \operatorname{Vol}(U_i)/\operatorname{Vol}(M_i)=O(\delta_i^2)\to0.
\]

\begin{proof}
On \(U_i\), \((|W_{+,i}|-\tfrac{1}{2\sqrt6})^2
\ge(\tfrac{2}{\sqrt6}-\tfrac{1}{2\sqrt6})^2
=\tfrac{9}{24}=\tfrac38\).
Hence \(\tfrac38\operatorname{Vol}(U_i)\le Y_i=O(A\delta_i^2)\),
giving \(\operatorname{Vol}(U_i)\le(8/3)O(A\delta_i^2)\),
and the volume fraction against \(\operatorname{Vol}(M_i)=O(A)\) is \(O(\delta_i^2)\).
\end{proof}

\noindent\textbf{Remark.}  This rules out more than the scalar-flat connector
in the familiar integral-minimizing constructions of Ishida--LeBrun~2003:
the negative K\"ahler--Einstein orbifold \emph{bulk} is also incompatible
with the Weyl profile.  For any region converging to a K\"ahler metric of
any scalar curvature \(s_0<0\), the Weyl magnitude is
\(|W_+|_{\rm K\"ahler}=|s_0|/(2\sqrt6)\), and the discrepancy from the
calibrated value \(1/(2\sqrt6)\) is \((|s_0|-1)/(2\sqrt6)\).
With \(s_0\to-4\), this is \(3/(2\sqrt6)\), giving the above exclusion.

Similarly, any region converging to \emph{real hyperbolic geometry}
(constant curvature \(-1/3\), \(W_+=0\)) has discrepancy
\(1/(2\sqrt6)\), hence occupies volume fraction at most \(O(A\delta_i^2\cdot24)
=O(A\delta_i^2)\) -- also vanishing.

\section{Bubble restriction to locally hyperk\"ahler}

\noindent\textbf{Corollary B2} (bubble restriction).
\emph{Suppose \(A\) is fixed (fixed manifold in the critical family).
Let \(U_i\subset M_i\) with \(\operatorname{Vol}(U_i)\to0\).
Then the self-dual Weyl energy vanishes:
\(\int_{U_i}|W_{+,i}|^2\,d\mu_i\to0\).}

\emph{Consequently, every smooth noncollapsed rescaling limit arising from
a vanishing-scale sequence is both Ricci-flat and self-dually flat
(\(W_+=0\)), hence locally hyperk\"ahler.}

\begin{proof}
\[
  \int_{U_i}|W_{+,i}|^2
  \le 2\int_{U_i}\bigl(|W_{+,i}|-\tfrac{1}{2\sqrt6}\bigr)^2
  +\tfrac{1}{12}\operatorname{Vol}(U_i)
  \le 2Y_i+\tfrac{1}{12}\operatorname{Vol}(U_i).
\]
With fixed \(A\): \(Y_i=O(A\delta_i^2)\to0\) and
\(\operatorname{Vol}(U_i)\to0\).

For a blow-up at scale \(\ell_i\to0\): rescale \(\hat g_i=\ell_i^{-2}g_i\).
The Weyl tensor is conformally invariant (as endomorphism of \(\Lambda^2_+\)),
hence \(|W_+(\hat g_i)|^2=\ell_i^4|W_+(g_i)|^2\) pointwise.
But the volume form scales as \(d\mu_{\hat g_i}=\ell_i^{-4}d\mu_{g_i}\),
so
\[
  \int_{B_{\hat g_i}(p,R)}|W_+(\hat g_i)|^2\,d\mu_{\hat g_i}
  =\int_{B_{g_i}(p,\ell_iR)}|W_+(g_i)|^2\,d\mu_{g_i}\to0
\]
since \(\operatorname{Vol}_{g_i}(B_{g_i}(p,\ell_iR))\to0\).  A smooth
convergent limit has \(W_+=0\) on every ball, hence globally.  Combined with
\(\operatorname{Ric}(\hat g_i)\to0\) (from the Step~3 scaling argument), the
limit is Ricci-flat and self-dually flat.  In four real dimensions, this forces
holonomy \(\subseteq\mathrm{SU}(2)\), i.e.\ locally hyperk\"ahler.
\end{proof}

\noindent\textbf{Remark} (varying topology).
If \(A_i\to\infty\), the bound becomes
\(\int_{U_i}|W_+|^2\le2A_i\delta_i^2+\operatorname{Vol}(U_i)/12\).
Absolute vanishing requires \(A_i\delta_i^2\to0\).  If
\(\delta_i\sim A_i^{-1/2}\), a bubble might carry finite \(W_+\) energy;
the restriction to hyperk\"ahler then fails and more exotic bubbles could
appear.  This quantifier distinction separates fixed-family rigidity from the
uniform statement needed for the original problem.

\section{Near-equality eigenvalue structure and the Armstrong tension}

The near-equality conditions from LeBrun~2006 Proposition~3.8 impose:
in the exact equality case, \(s-\sqrt6|W_+|=\text{const}<0\), the SW spinor
\(\Phi\) has \(|\Phi|^2=\text{const}\), and the self-dual 2-form
\(\sigma(\Phi)\) defines an almost-K\"ahler form \(\omega\) in the lowest
eigenspace of \(W_+\), with the two largest \(W_+\) eigenvalues equal.

At the calibrated values \(s=-4\), \(|W_+|=1/(2\sqrt6)\):
\begin{align}
  s-\sqrt6|W_+|&=-4-\tfrac12=-\tfrac92, \tag{E1}\\
  |W_+|^2&=\tfrac{1}{24},\quad
  \text{eigenvalues}\; (\tfrac{1}{12},\tfrac{1}{12},-\tfrac{1}{6}). \tag{E2}
\end{align}
The star-scalar curvature is
\[
  s^*=\tfrac{s}{3}+2W_+(\omega,\omega)
  =\tfrac{s}{3}-2\sqrt{\tfrac{3}{2}}|W_+|
  =-\tfrac43-\sqrt{\tfrac{3}{2}}\cdot\tfrac{1}{2\sqrt6}
\]
or equivalently from
\[
  s+s^*=\tfrac{4}{3}(s-\sqrt6|W_+|)=\tfrac43\cdot(-\tfrac92)=-6,
\]
giving \(s^*=-6-(-4)=-2\ne-4=s\).
For a K\"ahler metric, \(s^*=s\); hence this saturated almost-K\"ahler metric
is \emph{strictly almost-K\"ahler} (\(J\) not integrable).

\medskip
\noindent\textbf{Corollary C} (Armstrong--LeBrun tension).
\emph{If the near-equality sequence converged smoothly to a limit metric
\(g_\infty\) on any fixed critical manifold, then:}
\begin{enumerate}
\item \(g_\infty\) is Einstein (\(\operatorname{Ric}=-g\)) by smooth
  convergence of Ricci.
\item \(g_\infty\) saturates the scalar--Weyl estimate (\(K\to A\)
  reaches equality).
\item By Proposition~3.8, \(g_\infty\) is a saturated almost-K\"ahler metric.
\item By Corollary~2.5 (Armstrong's theorem),
  an Einstein saturated almost-K\"ahler metric must be K\"ahler--Einstein.
\item But \(|W_+|^2(g_\infty)=1/24\), whereas any K\"ahler metric with
  \(s=-4\) has \(|W_+|^2=s^2/24=2/3\).  Contradiction.
\end{enumerate}
\emph{Therefore, no sequence with \(\delta_i\to0\) on a fixed critical manifold
can converge smoothly to a single limit.}

\noindent\textbf{Remark.}  This is not yet a complete theorem excluding
\(\delta_i\to0\): we have not established the smooth convergence
hypothesis.  What the argument shows is that \emph{any mechanism allowing
\(\delta_i\to0\) must involve failure of smooth compactness}: either
collapsing, curvature concentration, or bubbling.  The bubble restriction
(Corollary~B2) limits the allowed degenerations to hyperk\"ahler ALE models
(Eguchi--Hanson, multi-Eguchi--Hanson, or more general crepant-resolution
models), but does not exclude all possibilities.

\section{Implication for the integral extremizing constructions}

The Ishida--LeBrun~2003 extremizers use three geometric pieces:
\begin{enumerate}
\item Negative K\"ahler--Einstein orbifold bulk (carrying most volume):
  \(|W_+|_{\rm KE}=2/\sqrt6\), discrepancy \(3/(2\sqrt6)\) from profile.
  \textbf{Ruled out by Corollary~B1}: volume fraction \(O(\delta^2)\).
\item Scalar-flat anti-self-dual connectors:
  \(|s+4|=4\), \(\delta\ge1\).
  \textbf{Already ruled out} (Step~3).
\item Ricci-flat gravitational instantons (ALE caps):
  these are Ricci-flat and can have \(W_+=0\)
  (hyperk\"ahler) or \(W_+\ne0\).  The profile allows
  only hyperk\"ahler caps with \(W_+=0\).
\end{enumerate}

None of these pieces can serve as the \emph{bulk} geometry of a
hypothetical \(\delta\to0\) sequence.  A sequence achieving the calibrated
Weyl profile must be supported by a fundamentally different geometric
construction --- one that is neither K\"ahler--Einstein, real hyperbolic,
conformally flat, nor anti-self-dual, but rather occupies the specific
intermediate state with \(W_+\) eigenvalues \((1/12,1/12,-1/6)\).

\section{Maximum-principle control on Seiberg--Witten data}

The Weitzenb\"ock formula
\[
  0=2\Delta|\Phi|^2+4|\nabla_A\Phi|^2+s|\Phi|^2+|\Phi|^4
\]
and the maximum principle give the pointwise bound
\[
  |\Phi|^2\le\max(-s,0)\le4(1+\delta_i)\quad\text{on }M_i.
\]
Hence \(|F_A^+|=|\Phi|^4/8\le(4+4\delta_i)^2/8\).

Both \(|\Phi|\) and \(|F_A^+|\) are uniformly bounded by the Ricci
pinching alone, \emph{without} any sectional curvature or harmonic-radius
hypothesis.  The near-equality further forces
\(|\Phi|^2\to|s|=4\) in \(L^6\) (from Proposition~3.8's H\"older
chain).  The remaining gap is between this \(L^6\) convergence and
pointwise \emph{non-vanishing} of \(|\Phi|\), which would give a
global almost-K\"ahler form.

\section{Summary and open quantitative problems}

The Weyl profile concentration (A1) provides a new geometric filter
on any hypothetical \(\delta\to0\) sequence:
\begin{itemize}
\item Known integral extremizers fail the filter (both KE bulk and
  scalar-flat connectors violate the profile).
\item Smooth noncollapsed blow-ups are restricted to hyperk\"ahler ALE
  (for fixed topology).
\item Smooth global convergence to a single limit is contradicted
  by Armstrong's theorem.
\end{itemize}

The precise open problems for completing the analysis are:
\begin{enumerate}
\item \textbf{Sobolev/harmonic-radius from Ricci pinching.}
  Can the uniform Ricci bound \(|\operatorname{Ric}+g|\le\delta g\)
  (with \(\delta\) small) together with the calibrated volume \(O(A)\)
  give a uniform Sobolev constant or harmonic-radius lower bound?
  If so, the smooth-convergence hypothesis in Corollary~C becomes
  automatic.
\item \textbf{Upgrade of \(|\Phi|^2\) convergence.}
  The \(L^6\) near-constancy of \(|\Phi|^2\) and the uniform \(L^\infty\)
  bound together require controlling the gradient of \(|\Phi|^2\).  The
  Weitzenb\"ock identity provides an equation for \(|\Phi|^2\), and
  Moser iteration or De~Giorgi--Nash--Moser estimates might yield
  H\"older continuity under the Ricci pinching alone.
\item \textbf{Non-vanishing of the almost-K\"ahler form.}
  If \(|\Phi|\to\text{const}\ne0\) pointwise, the self-dual form
  \(\sigma(\Phi)\) gives an almost-K\"ahler form with the forbidden
  Chern class on the critical manifold.  The zero set of \(|\Phi|\) is
  the obstruction; its codimension and topological structure are unknown
  under the present hypotheses.
\item \textbf{Phase separation vs.\ collapse topology.}
  If the almost-K\"ahler form cannot be made globally nonvanishing,
  the canonical classes must phase-separate across macroscopic regions,
  with transitions through thin topology cores.  Characterize these
  cores and determine whether the required topology can be carried
  through them while maintaining the Weyl profile.
\end{enumerate}

\medskip
\noindent\textbf{Files.}
Verification script:
\texttt{workspace/scripts/verify\_weyl\_profile.py}.
Step~3 base note:
\texttt{workspace/notes/step3\_boundary\_near\_saturation.tex}.
Step~3 constants:
\texttt{workspace/scripts/verify\_step3\_boundary\_constants.py} (passed).
